% Part: first-order-logic
% Chapter: tableaux
% Section: rules-and-proofs

\documentclass[../../../include/open-logic-section]{subfiles}

\begin{document}

\iftag{FOL}
      {\olfileid{fol}{tab}{rul}}
      {\olfileid{pl}{tab}{rul}}

\olsection{નિયમો અને \usetoken{P}{tableau}}

!!^a{tableau} એ !!a{structure}માં !!a{sentence} કઈ સંભવિત રીતોએ સત્ય
કે અસત્ય હોઈ શકે તેનો વ્યવસ્થિત સર્વે છે. ટેબ્લોના રચનાત્મક ઘટકો
!!{signed formula}s છે: !!{sentence}s સાથે સત્ય-મૂલ્યનું ``ચિહ્ન,'' જે
$\True$ અથવા~$\False$ હોય છે. આ ચિહ્નિત !!{formula}sને (નીચેની તરફ
વધતા) વૃક્ષમાં ગોઠવવામાં આવે છે.

\begin{defn}
  \emph{!!{signed formula}} એ સત્ય-મૂલ્ય અને !!a{sentence}થી બનેલી
  જોડી છે, એટલે કે આ બેમાંથી એક:
  \[
  \sFmla{\True}{!A} \text{ or } \sFmla{\False}{!A}.
  \]
\end{defn}

સહજ અર્થમાં, આપણે $\sFmla{\True}{!A}$ને ``$!A$ સત્ય હોઈ શકે'' અને
$\sFmla{\False}{!A}$ને ``$!A$ અસત્ય હોઈ શકે'' (કોઈ !!{structure}માં)
એમ વાંચી શકીએ.

વૃક્ષમાંનું દરેક !!{signed formula} કાં તો \emph{ધારણા} છે (ધારણાઓ
વૃક્ષની સાવ ટોચે લખાય છે), અથવા તેની ઉપરના !!a{signed formula}માંથી
કોઈ એક અનુમાન-નિયમ વડે મેળવવામાં આવ્યું છે. અગાઉના !!{formula}ના દરેક
સંભવિત !!{main operator} માટે બે નિયમ છે: ચિહ્ન~$\True$ હોય તે કિસ્સા
માટે એક અને ચિહ્ન~$\False$ હોય તે કિસ્સા માટે એક. કેટલાક નિયમો વૃક્ષની
શાખાઓ પાડવાની છૂટ આપે છે, જ્યારે કેટલાક શાખામાં માત્ર !!{signed formula}s
ઉમેરે છે. નિયમને તેની તુરંત અગાઉના !!{signed formula}ને જ લાગુ કરવો
જરૂરી નથી; તેને મૂળથી નિયમ લાગુ કરવાના સ્થાન સુધીની શાખામાંના કોઈપણ
!!{signed formula}ને લાગુ કરી શકાય છે (અને ઘણી વાર તેમ કરવું પડે છે).

કોઈ શાખામાં $\sFmla{\True}{!A}$ અને $\sFmla{\False}{!A}$ બંને હોય ત્યારે
તે \emph{બંધ} કહેવાય છે. જે ટેબ્લોની દરેક શાખા બંધ હોય તે બંધ
!!{tableau} છે. સહજ અર્થઘટન મુજબ, દરેક શાખા એક સંયુક્ત સંભાવના દર્શાવે
છે, પરંતુ $\sFmla{\True}{!A}$ અને $\sFmla{\False}{!A}$ સંયુક્ત રીતે
સંભવિત નથી. બીજા શબ્દોમાં, શાખા બંધ હોય તો તે જે સંભાવના દર્શાવે છે તે
નકારી કાઢવામાં આવી છે. ખાસ કરીને, તેનો અર્થ એ થાય છે કે બંધ !!{tableau}
$\sFmla{\True}{!A}$ સ્વરૂપની દરેક ધારણાને એકસાથે સત્ય અને
$\sFmla{\False}{!A}$ સ્વરૂપની દરેક ધારણાને એકસાથે અસત્ય બનાવવાની બધી
સંભાવનાઓ નકારી કાઢે છે.

$!A$ \emph{માટેનો} બંધ !!{tableau} એ એવો બંધ !!{tableau} છે જેનું
મૂળ~$\sFmla{\False}{!A}$ છે. આવો બંધ !!{tableau} અસ્તિત્વમાં હોય તો
$!A$ અસત્ય હોય તેવી બધી સંભાવનાઓ નકારી કાઢવામાં આવી છે; એટલે કે $!A$
દરેક !!{structure}માં સત્ય જ હોવું જોઈએ.

\end{document}
